\chapter*{Resumen} % OBLIGATORIO
RutaCuatri es un planificador de la trayectoria académica de un estudiante irregular en un sistema cuatrimestral. Dado un plan de estudios, el historial del estudiante y una meta (toda la carrera, llegar a la estadía o una especialidad), encuentra qué materias inscribir en cada cuatrimestre con el menor costo, medido como el número de cuatrimestres más una penalización por cada crédito arriba de la carga normal, y respetando la seriación, la oferta de cada materia por periodo y la carga máxima.

El problema se modela como una búsqueda en el espacio de pares (materias aprobadas, periodo), que tiene hasta $3 \cdot 2^{N}$ estados: más de 805 millones con 28 materias. Se resuelve con Anytime A*, que ejecuta Weighted A* con pesos 3, 2, 1.5 y 1: entrega un primer plan casi de inmediato, lo mejora en cada pasada con la garantía de costar como mucho $H$ veces el óptimo, y la última pasada demuestra el óptimo. La heurística $h = \max(\lceil R/C_{max} \rceil, L)$ combina los créditos que faltan con la cadena de prerrequisitos más larga y es admisible y consistente.

El sistema es una aplicación web con el motor de búsqueda separado de la interfaz y ejecutado en un Web Worker. Se probó con tres mallas de 28 materias, doce estudiantes tipo y nueve metas: en los 77 casos de prueba coincide con una implementación de referencia, y en 52 de 72 escenarios la primera pasada ya da el óptimo. Los experimentos muestran que las reglas de poda reducen las expansiones hasta 19 veces sin cambiar el óptimo y que la carga normal es el parámetro que más decide cuándo egresa un estudiante.
